\section{Pengujian}
% Nama berkas dan perintah panjang di \texttt tidak bisa dipenggal.
\sloppy

Seluruh pengujian pada bab ini berjalan otomatis dengan \texttt{mix test}, kecuali skenario
C01 sampai C05 yang berjalan pada pipeline anmategra. \textit{Workflow} \texttt{ci.yml}
menjalankan tes dua kali: pada runner Ubuntu, dan di dalam image
\texttt{elixir:1.20-otp-29-alpine} yang memakai musl dan OpenSSL yang sama dengan image
rilis. Image Alpine tersebut tidak memuat CLI \texttt{openssl}, sehingga tes
interoperabilitas hanya berjalan pada runner Ubuntu.

Untuk laporan ini, seluruh tes dijalankan ulang di dalam image Alpine yang sama dengan
CLI OpenSSL 3.5.9 ditambahkan, sehingga tes interoperabilitas ikut berjalan. Hasilnya 213
tes lulus (180 tes dan 33 doctest) tanpa kegagalan dan tanpa tes yang dilewati. Property
test pada \texttt{test/property/} dijalankan dengan cara yang sama dan 20 property lulus.

\subsection{Skenario U01 sampai U14}

Skenario ditulis pada berkas \texttt{test/cli\_test.exs} dan memanggil CLI lewat
\texttt{Nadina.CLI.run/1}, fungsi yang sama dengan titik masuk \texttt{bin/nadina} tanpa
\texttt{halt}. Exit code dan keluaran \texttt{verify -json} diperiksa langsung. Berkas
masukan dibentuk oleh tes di direktori sementara (\texttt{@tag :tmp\_dir}) dengan nilai
dasar berikut:

\begin{itemize}
  \item artefak \texttt{contoh-app.zip} berisi teks \texttt{isi artefak},
  \item paket \texttt{contoh-app} dan versi \texttt{1.0.0},
  \item satu pasangan kunci per profil hasil \texttt{nadina keygen},
  \item \textit{trust store} hasil \texttt{nadina trust add} yang mengizinkan kunci
    tersebut untuk \texttt{contoh-app},
  \item bundle hasil \texttt{nadina sign} atas artefak di atas.
\end{itemize}

U01 dan U14 dijalankan untuk ketiga profil. U02 sampai U13 memakai profil
\texttt{ecdsa-p256-sha256}, karena pemeriksaannya tidak bergantung pada algoritma.
Setiap tes mengubah satu hal dari kondisi dasar, lalu memverifikasi ulang.

{\small
\begin{longtable}{@{}p{1cm}p{7.2cm}p{4.3cm}p{1.6cm}@{}}
\toprule
ID & Perubahan pada masukan & Hasil yang diharapkan & Hasil \\
\midrule
\endhead
U01 & Berkas teks (\texttt{halo dunia}), berkas biner 5000 byte acak, dan berkas kosong
  ditandatangani lalu diverifikasi, untuk tiap profil. Hasil juga diperiksa memuat
  \texttt{algorithm}, \texttt{keyid}, \texttt{package}, dan \texttt{version} yang benar.
  & \texttt{OK}, exit 0 & Lulus \\
U02 & Satu byte artefak diganti (\texttt{isi artefak} menjadi \texttt{isi artefaK}),
  ukuran dan nama tetap. & \texttt{ARTIFACT\_MISMATCH}, exit 1 & Lulus \\
U03 & Satu byte ditambahkan di akhir artefak, lalu pada uji terpisah satu byte dihapus.
  & \texttt{ARTIFACT\_MISMATCH}, exit 1 & Lulus \\
U04 & Versi manifest diganti menjadi \texttt{9.9.9} dan diserialisasi ulang secara
  canonical. Tanda tangan lama dipertahankan dan verifikasi memakai
  \texttt{-version 9.9.9}. & \texttt{INVALID\_SIGNATURE}, exit 1 & Lulus \\
U05 & Bit terendah pada byte pertama tanda tangan dibalik. Panjang tanda tangan tetap dan
  Base64 tetap sah. & \texttt{INVALID\_SIGNATURE}, exit 1 & Lulus \\
U06 & Verifikasi dengan \textit{trust store} kosong, lalu dengan \textit{trust store} asli
  setelah kunci penerbit dihapus lewat \texttt{trust remove}.
  & \texttt{KEY\_NOT\_TRUSTED}, exit 1 & Lulus \\
U07 & \textit{Trust store} berisi kunci penerbit yang hanya diizinkan untuk paket
  \texttt{paket-lain}. & \texttt{KEY\_NOT\_TRUSTED}, exit 1 & Lulus \\
U08 & Bundle sah diverifikasi dengan \texttt{-version 2.0.0}.
  & \texttt{METADATA\_MISMATCH}, exit 1 & Lulus \\
U09 & Artefak disalin ke \texttt{nama-lain.zip} tanpa mengubah isi, lalu berkas salinan
  yang diverifikasi. & \texttt{METADATA\_MISMATCH}, exit 1 & Lulus \\
U10 & Empat bundle diuji terpisah: field \texttt{manifest\_b64} ganda,
  \texttt{manifest\_b64} bertipe integer, Base64 tanda tangan tidak sah, dan manifest
  sebesar 16 KiB ditambah 1 byte. & \texttt{FORMAT\_ERROR}, exit 1, \texttt{algorithm}
  bernilai \texttt{null} pada tiap kondisi & Lulus \\
U11 & Nilai \texttt{algorithm} pada manifest diganti \texttt{rsa-sha1}.
  & \texttt{UNSUPPORTED\_ALGORITHM}, exit 1 & Lulus \\
U12 & Artefak ditandatangani dengan kunci ECDSA P-256 lain yang tidak ada di
  \textit{trust store}. \texttt{keyid} pada manifest sesuai kunci tersebut.
  & \texttt{KEY\_NOT\_TRUSTED}, exit 1, \texttt{keyid} pada hasil sama dengan fingerprint
  kunci baru & Lulus \\
U13 & Bundle ditulis ulang dengan spasi, tab, dan baris baru tambahan serta urutan field
  dibalik. Kedua string Base64 tidak diubah. & \texttt{OK}, exit 0 & Lulus \\
U14 & Manifest identik ditandatangani dua kali dengan kunci yang sama, untuk tiap profil.
  & Kedua bundle \texttt{OK}. Tanda tangan RSA-PSS dan ECDSA berbeda, Ed25519 identik.
  & Lulus \\
\bottomrule
\caption{Skenario pengujian U01 sampai U14.}
\end{longtable}
}

Hasil U14 sesuai dengan sumber randomness tiap algoritma. RSA-PSS memakai salt 32 byte
acak dan ECDSA memakai nonce $k$ acak, sehingga dua tanda tangan atas pesan yang sama
berbeda. Ed25519 menurunkan nonce dari kunci dan pesan (RFC 8032), sehingga tanda
tangannya identik. Analisis lengkapnya, termasuk 200 tanda tangan ECDSA tanpa nilai $r$
berulang, ada pada subbab Nonce dan Randomness di bab Analisis Keamanan.

Selain tabel di atas, \texttt{cli\_test.exs} menguji kasus tambahan pada urutan status
verifikasi: opsi wajib yang hilang dan berkas yang tidak ada menghasilkan
\texttt{INPUT\_ERROR} dengan exit 2, \textit{trust store} rusak atau fingerprint yang tidak
cocok dengan kuncinya menghasilkan \texttt{FORMAT\_ERROR}, dan kunci tepercaya dengan
profil berbeda dari manifest menghasilkan \texttt{KEY\_NOT\_TRUSTED}.

\subsubsection*{Cara menjalankan}
\begin{lstlisting}[language=bash,numbers=none]
mix test test/cli_test.exs
mix test test/cli_test.exs --only describe:"U14: dua tanda tangan atas manifest identik"
\end{lstlisting}

\subsection{Test Vector dan Property Test}

Setiap profil diuji terhadap vektor dari sumber publik, terpisah dari CLI:

\begin{itemize}
  \item RSA-PSS: 108 kasus Wycheproof \texttt{rsa\_pss\_3072\_sha256\_mgf1\_32}, yang
    parameternya sama persis dengan profil (RSA 3072, SHA-256, MGF1-SHA-256, salt 32 byte).
    Semua kasus \textit{valid} diterima dan semua kasus \textit{invalid} ditolak.
  \item ECDSA P-256: vektor RFC 6979 lampiran A.2.5 untuk pesan \texttt{sample}. Tanda
    tangan diterima, pesan \texttt{Sample} ditolak, dan kunci privat vektor menurunkan
    titik publik yang sama.
  \item Ed25519: RFC 8032 bagian 7.1 TEST 1 (pesan kosong) dan TEST 2 (pesan 1 byte).
    Kunci publik, tanda tangan, dan hasil verifikasi sama dengan vektor.
\end{itemize}

Tes profil juga memeriksa nilai di luar domain tanpa \textit{raise}: $r$ atau $s$ bernilai
$0$, $n$, atau $2^{256}-1$ pada ECDSA, serta tanda tangan bernilai $0$, $1$, $n-1$, $n$,
dan $2^{3072}-1$ pada RSA-PSS. Semuanya menghasilkan \texttt{false}.

Property test di \texttt{test/property/} memakai StreamData untuk membangkitkan masukan
acak:

\begin{itemize}
  \item \texttt{strict\_json}: dokumen JSON sah diterima dengan nilai yang sama, sedangkan
    \textit{duplicate key}, \texttt{null} di posisi mana pun, data setelah nilai pertama,
    dan boolean di posisi integer ditolak.
  \item manifest dan bundle: \texttt{parse(serialize(m))} mengembalikan manifest dan byte
    yang sama. White space, urutan field selain canonical, field tambahan, serta boolean
    atau pecahan di posisi \texttt{size} ditolak. Pada bundle, white space dan urutan field
    tidak mengubah isi, sedangkan Base64 tanpa padding atau dengan bit padding tidak nol
    ditolak.
  \item sign dan verify: untuk tiap profil, 40 rilis acak yang ditandatangani lewat CLI
    selalu menghasilkan \texttt{OK}, dan 50 pesan acak menghasilkan tanda tangan sah yang
    ditolak untuk pesan lain.
\end{itemize}

\subsubsection*{Cara menjalankan}
\begin{lstlisting}[language=bash,numbers=none]
mix test test/profile test/property test/nonce_analysis_test.exs
\end{lstlisting}

\subsection{Interoperabilitas}

Interoperabilitas diuji terhadap OpenSSL untuk ketiga profil (berkas
\texttt{test/interop/}). OpenSSL hanya dipakai sebagai alat uji, tidak oleh program. Tes
bertag \texttt{:interop} dilewati bila \texttt{openssl} tidak ada di \texttt{PATH}.

\begin{table}[H]
\centering
\small
\begin{tabular}{@{}p{2.6cm}p{4cm}p{4cm}p{3.4cm}@{}}
\toprule
Profil & Nadina ke OpenSSL & OpenSSL ke Nadina & Harus ditolak \\
\midrule
RSA-PSS &
  Kunci dibaca OpenSSL (3072 bit, $e = 65537$), fingerprint sama. Tanda tangan diterima
  \texttt{openssl pkeyutl -verify} dengan opsi PSS salt 32 dan MGF1-SHA-256. &
  Kunci buatan OpenSSL dimuat sebagai profil RSA-PSS dan tanda tangannya sah. &
  PSS dengan salt lain dan PKCS\#1 v1.5 dari OpenSSL. \\
\addlinespace
ECDSA P-256 &
  Kurva dikenali sebagai \texttt{prime256v1}, fingerprint sama. Tanda tangan $r\|s$ diubah
  ke DER lalu diterima \texttt{openssl dgst -verify}. &
  DER dari OpenSSL diubah ke $r\|s$ lalu diterima, termasuk banyak tanda tangan dengan
  panjang DER berbeda-beda. &
  Pesan yang diubah. \\
\addlinespace
Ed25519 &
  Kurva dikenali sebagai Ed25519. SPKI dari OpenSSL tanpa field parameter dan fingerprint
  sama. Tanda tangan diterima OpenSSL. &
  Kunci buatan OpenSSL dimuat sebagai Ed25519 dan tanda tangannya sah. &
  Pesan yang diubah dan tanda tangan atas hash pesan (pola Ed25519ph). \\
\bottomrule
\end{tabular}
\caption{Cakupan uji interoperabilitas dengan OpenSSL.}
\end{table}

Pada ketiga profil, kunci publik yang diturunkan OpenSSL dari kunci privat Nadina sama
dengan berkas \texttt{.pub.pem} yang ditulis Nadina, dan fingerprint
$\mathrm{SHA256}(\text{DER SPKI})$ dari kedua sisi sama.

\subsubsection*{Cara menjalankan}
\begin{lstlisting}[language=bash,numbers=none]
mix test --only interop
\end{lstlisting}

\subsection{Pengujian Integrasi CI/CD (C01 sampai C05)}
% TODO (Nayaka): tautan workflow run per skenario, referensi image dan digest, log, keluaran
% JSON verify, bukti ada atau tidak adanya berkas penanda. Bukti disimpan di docs/ci-evidence/.

\subsection{Pengujian Performa}

\subsubsection*{Lingkungan}
Benchmark dijalankan sekali pada satu mesin, di dalam image rilis Nadina (Docker Desktop
dengan backend WSL2) dan dengan \texttt{--network none}. Data mentah dan ringkasan ada di
\texttt{data/benchmark/}, dan tabel pada subbab ini dibaca langsung dari berkas CSV
tersebut.

Container diberi hostname \texttt{nadina} yang dipetakan ke \texttt{127.0.0.1} lewat
\texttt{--add-host}. Pada percobaan awal tanpa pemetaan ini, Docker Desktop tidak menuliskan
hostname container ke \texttt{/etc/hosts}, sehingga setiap proses BEAM menunggu timeout
DNS ke resolver yang tidak terjangkau. Akibatnya \texttt{nadina -version} memerlukan
5,55 detik, dibandingkan 0,47 detik setelah hostname dipetakan. Selisih ini berasal dari
lingkungan, bukan dari Nadina, tetapi dapat muncul pada pemakaian \texttt{--network none}
di lingkungan lain.

\begin{table}[H]
\centering
\small
\begin{tabular}{@{}ll@{}}
\toprule
Komponen & Nilai \\
\midrule
CPU & 13th Gen Intel Core i7-13700HX, 24 CPU logis \\
Memori untuk VM WSL2 & 7,6 GiB \\
OS host & Windows 11 Home Single Language 10.0.26200 \\
Kernel & 6.18.58.1-microsoft-standard-WSL2 (x86\_64) \\
OS container & Alpine Linux 3.24 (musl) \\
Runtime & Elixir 1.20.5, Erlang/OTP 29 (ERTS 17.1.1) \\
Pustaka kripto & \texttt{:crypto} dan \texttt{:public\_key} OTP di atas OpenSSL 3.5.9 \\
\bottomrule
\end{tabular}
\caption{Lingkungan benchmark (\texttt{data/benchmark/env.csv}).}
\end{table}

\subsubsection*{Metode}
Skrip \texttt{bench/bench.exs} dijalankan lewat \texttt{bench/run.sh}. Setiap operasi
dijalankan 5 kali sebagai pemanasan lalu diukur 30 kali, kecuali pembangkitan kunci yang
diukur 5 pasang per profil. Waktu diambil dari \texttt{:erlang.monotonic\_time/0}, sehingga
tidak terpengaruh perubahan jam sistem. Median memakai definisi biasa (rata-rata dua nilai
tengah bila $n$ genap). Persentil ke-95 memakai metode \textit{nearest-rank}: data diurutkan
naik dan P95 adalah elemen ke-$\lceil 0{,}95 \cdot n \rceil$ dengan posisi dimulai dari 1.
Untuk $n = 30$, P95 adalah elemen ke-29.

Artefak berukuran 64 KiB, 4 MiB, dan 64 MiB berisi byte acak, dan berkas yang sama dipakai
untuk ketiga profil. Komponen yang diukur sebagai berikut:

\begin{itemize}
  \item pembangkitan kunci: \texttt{Keys.generate/1}, terpisah dari sign dan verify,
  \item hash artefak: SHA-256 streaming atas byte artefak (\texttt{Artifact.digest/1}),
  \item sign dan verify: operasi kriptografi atas byte manifest $M$ di dalam proses yang
    sama, termasuk hash internal algoritma, tanpa I/O dan tanpa hash artefak,
  \item end-to-end: \texttt{nadina sign} dan \texttt{nadina verify -json} sebagai proses
    baru, termasuk boot BEAM, I/O, hash artefak, parsing, dan penulisan keluaran.
    \texttt{nadina -version} diukur dengan cara yang sama sebagai pembanding biaya boot.
\end{itemize}

\pgfplotstableset{
  col sep=comma,
  every head row/.style={before row=\toprule, after row=\midrule},
  every last row/.style={after row=\bottomrule},
  bench/.style={fixed, fixed zerofill, precision=1, use comma, 1000 sep={.}},
}

\subsubsection*{Pembangkitan kunci}

\begin{table}[H]
\centering
\pgfplotstabletypeset[bench,
  columns={algorithm,n,median_us,min_us,max_us},
  columns/algorithm/.style={string type, column name=Profil, column type=l},
  columns/n/.style={column name=$n$, int detect},
  columns/median_us/.style={column name=Median (ms), preproc/expr={##1/1000}, precision=2},
  columns/min_us/.style={column name=Min (ms), preproc/expr={##1/1000}, precision=2},
  columns/max_us/.style={column name=Maks (ms), preproc/expr={##1/1000}, precision=2},
]{data/benchmark/keygen.csv}
\caption{Waktu pembangkitan satu pasangan kunci.}
\end{table}

\subsubsection*{Hash artefak}

\begin{table}[H]
\centering
\pgfplotstabletypeset[bench,
  columns={size_label,n,median_us,p95_us,throughput_mib_s},
  columns/size_label/.style={string type, column name=Artefak, column type=l},
  columns/n/.style={column name=$n$, int detect},
  columns/median_us/.style={column name=Median (ms), preproc/expr={##1/1000}, precision=2},
  columns/p95_us/.style={column name=P95 (ms), preproc/expr={##1/1000}, precision=2},
  columns/throughput_mib_s/.style={column name=Throughput (MiB/s), precision=0},
]{data/benchmark/hash.csv}
\caption{Waktu SHA-256 atas artefak. Throughput dihitung dari median.}
\end{table}

\subsubsection*{Sign dan verify atas manifest}

\begin{table}[H]
\centering
\pgfplotstabletypeset[bench,
  columns={algorithm,op,n,median_us,p95_us},
  columns/algorithm/.style={string type, column name=Profil, column type=l},
  columns/op/.style={string type, column name=Operasi, column type=l},
  columns/n/.style={column name=$n$, int detect},
  columns/median_us/.style={column name=Median ($\mu$s)},
  columns/p95_us/.style={column name=P95 ($\mu$s)},
]{data/benchmark/sign_verify.csv}
\caption{Waktu operasi kriptografi atas byte manifest $M$.}
\end{table}

\begin{figure}[H]
\centering
\begin{tikzpicture}
\begin{axis}[
  ybar, bar width=14pt, width=0.8\textwidth, height=6cm,
  ymode=log, log origin=infty,
  ylabel={Median ($\mu$s, skala log)},
  symbolic x coords={rsa-pss-sha256,ecdsa-p256-sha256,ed25519},
  xtick=data, enlarge x limits=0.25,
  legend style={at={(0.5,-0.18)}, anchor=north, legend columns=-1},
]
\addplot table[x=algorithm, y=median_us, col sep=comma, only op=sign]
  {data/benchmark/sign_verify.csv};
\addplot table[x=algorithm, y=median_us, col sep=comma, only op=verify]
  {data/benchmark/sign_verify.csv};
\legend{sign, verify}
\end{axis}
\end{tikzpicture}
\caption{Median sign dan verify atas manifest per profil.}
\end{figure}

\subsubsection*{End-to-end}

\begin{table}[H]
\centering
\small
\pgfplotstabletypeset[bench,
  columns={algorithm,op,size_label,median_us,p95_us},
  columns/algorithm/.style={string type, column name=Profil, column type=l},
  columns/op/.style={string type, column name=Perintah, column type=l},
  columns/size_label/.style={string type, column name=Artefak, column type=l},
  columns/median_us/.style={column name=Median (ms), preproc/expr={##1/1000}, precision=0},
  columns/p95_us/.style={column name=P95 (ms), preproc/expr={##1/1000}, precision=0},
]{data/benchmark/e2e.csv}
\caption{Waktu end-to-end per perintah sebagai proses baru. Baris \texttt{boot} adalah
\texttt{nadina -version}.}
\end{table}

\subsubsection*{Ukuran}

\begin{table}[H]
\centering
\pgfplotstabletypeset[bench,
  columns={algorithm,private_der,public_der,signature,manifest,bundle},
  columns/algorithm/.style={string type, column name=Profil, column type=l},
  columns/private_der/.style={column name=Privat DER, int detect},
  columns/public_der/.style={column name=Publik DER, int detect},
  columns/signature/.style={column name=Tanda tangan, int detect},
  columns/manifest/.style={column name=Manifest, int detect},
  columns/bundle/.style={column name=Bundle, int detect},
]{data/benchmark/sizes.csv}
\caption{Ukuran dalam byte. Kunci privat dalam PKCS\#8 dan kunci publik dalam SPKI, keduanya
DER tanpa enkripsi. Tanda tangan diukur sebelum Base64.}
\end{table}

\subsubsection*{Analisis}
% TODO: isi setelah benchmark penuh (runs=30) selesai.

\subsection{Pengujian SBOM}
Pengujian SBOM berada di \texttt{test/sbom\_test.exs} dengan 16 tes, dan lolos seluruhnya.
Fixture berupa keluaran asli Trivy dalam format CycloneDX dan SPDX, yang disimpan di
\texttt{test/testdata/sbom/}. Alur ujung ke ujung juga dijalankan di dalam image dengan UID acak
dan \texttt{--network none}: verifikasi valid menghasilkan exit 0, artefak yang diubah
menghasilkan exit 1, dan hasil \texttt{extract} identik byte demi byte dengan keluaran Trivy.

\begin{table}[H]
\centering
\begin{tabular}{@{}lll@{}}
\toprule
Kasus uji & Hasil yang diharapkan & Hasil \\
\midrule
SBOM valid dan saling merujuk & OK, exit 0 & OK, exit 0 \\
CycloneDX dan SPDX dari Trivy di-bind lalu di-extract & byte identik & byte identik \\
Masukan SBOM bukan objek JSON & ditolak saat bind, exit 2 & ditolak, exit 2 \\
\textit{Wrapper} diubah & ARTIFACT\_MISMATCH, exit 1 & ARTIFACT\_MISMATCH, exit 1 \\
Artefak rilis diganti isinya & ARTIFACT\_MISMATCH, exit 1 & ARTIFACT\_MISMATCH, exit 1 \\
\textit{Wrapper} milik artefak lain, bundle sah & ARTIFACT\_MISMATCH, exit 1 & ARTIFACT\_MISMATCH, exit 1 \\
Bundle \textit{wrapper} dari kunci tak tepercaya & KEY\_NOT\_TRUSTED, exit 1 & KEY\_NOT\_TRUSTED, exit 1 \\
\textit{Wrapper} rusak namun ditandatangani & FORMAT\_ERROR, exit 1 & FORMAT\_ERROR, exit 1 \\
Isi SBOM bukan JSON namun \textit{wrapper} sah & FORMAT\_ERROR, exit 1 & FORMAT\_ERROR, exit 1 \\
Berkas \textit{wrapper} tidak ada & INPUT\_ERROR, exit 2 & INPUT\_ERROR, exit 2 \\
\bottomrule
\end{tabular}
\caption{Hasil pengujian SBOM.}
\end{table}

\fussy
